\subsection{The all-type and specified sequential path comparisons}
Throughout this subsection categories may be arbitrary small categories.
A marking is a subcategory containing all objects; it need not contain
all isomorphisms. Fix endpoints $a,b$ in $(M,\Sigma)$. Write
$\mathcal A(M;a,b)$ for the fibred all-type category: an object is
$p:I_\tau\to M$ with marked negative edges, and an arrow from $p$ to
$q:I_\sigma\to M$ is a type map $u:I_\tau\to I_\sigma$ and an
endpoint-fixed transformation $p\Rightarrow qu$. Write
$\mathcal W(M;a,b)$ for the same category with every component of this
transformation required to lie in $\Sigma$.

For an actual projective type sequence $t:\N^{\op}\to\mathrm{Ty}$,
write $\mathcal E_t(M;a,b)$ for its fibred total category. Thus an
arrow from stage $n$ to stage $m$ requires $n\geq m$ and is a
transformation $p\Rightarrow q\,t(n\to m)$. These conventions specify
the functors, not just the spaces up to an unspecified equivalence.

\begin{lemma}[The marked comparison and its marking hypothesis]
\label{F:marked-hammocks}
There is a natural weak equivalence, expressed by the canonical
Grothendieck and reduction zigzag,
\[
 B\mathcal W(M;a,b)\simeq
 \operatorname{Map}_{L(M,\Sigma)}(a,b).
\]
It respects concatenation. This statement holds for the marking
convention just specified.
\end{lemma}
\begin{proof}
The fixed-type nerves are the fixed-endpoint marked hammock diagrams.
The type category has its Reedy factorization into a vertex-surjection
followed by a vertex-injection. In the opposite index, a latching map
is the inclusion of diagrams with the relevant collapsed columns equal
to identities; intersections correspond to the generated common
quotient. Its latching maps are consequently monomorphisms. The
matching indexing categories for constant diagrams are empty or
connected: the reduced interval consisting of the endpoints and the
direction-change vertices is initial among proper injective
subintervals whenever there is one. Hence ordinary colimit computes
homotopy colimit. These are the precise combinatorial checks in
\cite[Proposition~4.8, Lemma~4.16]{arakawacnossen2026}; they use
identities and closure under composition, not inclusion of every
isomorphism.

The colimit is the reduced hammock simplicial set: identity columns
are deleted and adjacent equally directed columns are composed.
\cite[Section~2.1 and Proposition~2.2]{dwyerkan1980calculating}
identify this with the simplicial localization. The convention is
explicit in \cite[Section~3.1]{dwyerkan1980function}: the subcategory
contains all objects. Thus no enlargement of $\Sigma$ has occurred.
For the infinity-categorical universal property, one may now enlarge
to the closure $\overline\Sigma$ of \cite[Section~3.4 and
Lemma~5.1]{dwyerkan1980simplicial}; that lemma makes the induced
simplicial-localization map a weak equivalence. This closure contains
every isomorphism, so \cite[Theorem~4.7 and Corollary~1.11]{arakawacnossen2026}
applies with its stated convention. A functor to any infinity-category
which inverts $\Sigma$ also inverts $\overline\Sigma$: pass to its
homotopy category and use the defining ordinary localization.
Consequently the two infinity-localization universal properties agree.
This explains, rather than assumes, the passage between conventions.

For orientation, if $F:T^{\op}\to\Cat$, then
\[
 \left(\int_T^{\mathrm{fib}}F\right)^{\op}
   =\int_{T^{\op}}^{\mathrm{opf}}F^{\op}.
\]
Reversing a nerve simplex gives a natural homeomorphism
$B(C^{\op})\cong BC$. Apply the verified Grothendieck homotopy-colimit
theorem to the right side and these homeomorphisms to the fibres.
This is the required comparison for the fibred convention used here.
All reduction operations and these natural comparisons commute with
block concatenation, which gives the last assertion.
\end{proof}

\begin{lemma}[A relative interval contracts all arbitrary-vertical paths]
\label{F:relative-interval}
Give $I_\tau$ the marking generated by its negative edges. Every
nonempty category $F_\rho(0,|\tau|)$ of relative paths in this interval,
with arbitrary endpoint-fixed transformations, is categorically
contractible. It is nonempty exactly when there is a type map
$\rho\to\tau$. Consequently $B\mathcal A(I_\tau)$ is contractible.
If $t$ is universally cofinal in the reachability sense of C.20, then
$B\mathcal E_t(I_\tau)$ is contractible as well.
\end{lemma}
\begin{proof}
For such a path $h$, regarded as a sequence of target vertices, set
\[
 \overline h(k)=\max_{j\leq k}h(j).
\]
A path can cross a positive target edge backwards neither on a
positive source edge (there is no such target arrow) nor on a negative
source edge (its image would be unmarked). Every numerical decrease
therefore stays within a negative run. The segment
$\overline h(k)\to h(k)$ is a marked arrow. Moreover
$\overline h$ is nondecreasing and preserves the required arrow
directions: when its value increases, its new segment is a subsegment
of the arrow traversed by the corresponding edge of $h$. In particular
a negative source edge still has marked image. Endpoints are fixed.
Thus $\overline h$ is a type map and there is an endpoint-fixed
transformation $\overline h\Rightarrow h$.

Numerical maximum is order-preserving separately in each variable on
any oriented-line poset: clipping a directed segment at a fixed
vertex gives a subsegment or an identity. Hence $h\mapsto\overline h$
is a functor on the path category, and the preceding transformation
is natural in $h$. It retracts onto $P(\rho,\tau)$ and is the identity
there. Lemma~\ref{F:type-contraction} proves categorical
contractibility and the nonemptiness assertion.

The supported full subcategory of types is contractible. The
endofunctor $\rho\mapsto\rho\tau$ has natural type maps to both
$\rho$ and the constant type $\tau$, obtained by collapsing the
other block. Its values are supported, and $\tau$ itself is supported.
The Grothendieck theorem, with the orientation just checked, now
applies to the contractible nonempty fibres. For a universal sequence,
the supported stages are exactly a nonempty terminal tail of $\N$;
the same argument applies. Empty target type causes no exception:
every relative path in its one-object identity category is unique.
The independent program \texttt{a\_f\_interval\_prefixmax.py} checks
all $961$ pairs of types of length at most four: $1756$ relative
endpoint functors, including $285$ nonmonotone ones, $9219$ arbitrary
transformations and $7208$ marked transformations. All stated
functoriality and marked-component assertions pass. A further $4565$
checks verify compatibility with source type pullback in lengths at
most three. These finite checks corroborate the general segment proof.
\end{proof}

\begin{theorem}[The missing C.3/C.7 comparisons, with specified maps]
\label{F:unmarked-path-comparison}
For every small $(M,\Sigma)$ and every pair of endpoints the inclusion
\[
 B\mathcal W(M;a,b)\longrightarrow B\mathcal A(M;a,b)
\]
is a weak equivalence. For every universally cofinal actual projective
sequence $t$, both maps in the specified span
\[
 B\colim_n F_{t_n}(a,b)\ \longleftarrow\
 B\mathcal E_t(M;a,b)\ \longrightarrow\
 B\mathcal A(M;a,b)
\]
are weak equivalences. These give the asserted localization mapping
space and its compatibility with composition. In particular the
path-model prerequisites of the supplied C.28 are valid with arbitrary
interior vertical arrows, under the exact marking convention of C.20.
\end{theorem}
\begin{proof}
\textbf{A universal diagram for one simplex.}
Let $\sigma=(p_0\to\cdots\to p_r)$ be a simplex in
$N\mathcal A(M;a,b)$. Write its type maps as $u_{ij}:I_{\tau_i}\to
I_{\tau_j}$ and its transformations as $\eta_{ij}:p_i\Rightarrow
p_ju_{ij}$. Form the covariant Grothendieck construction of the
interval diagram over $[r]$. Its objects are $(i,k)$; an arrow to
$(j,l)$, for $i\leq j$, is an arrow $u_{ij}(k)\to l$ in
$I_{\tau_j}$. Composition applies $u_{jv}$ and then composes in
$I_{\tau_v}$. This formula proves the category laws.

Identify the left endpoint column with one object $A$ and the right
endpoint column with one object $B$, and set their vertical arrows to
identities. Take this quotient as a categorical presentation, denoted
$G_\sigma$. It is small, although the quotient can have infinitely
many arrows. Evaluation at the given diagram is a strict functor
$G_\sigma\to M$: the defining equations are precisely naturality
and composition for the $\eta_{ij}$.

Mark the negative row arrows and, for every pair $i\leq j$, every
cocartesian arrow $(i,k)\to(j,u_{ij}k)$ whose evaluated component
$\eta_{ij}(k)$ belongs to $\Sigma$; close under identities and
composition. This all-pairs condition is essential. It makes
$G_\theta\to G_\sigma$ relative for every face $\theta\subseteq
\sigma$, including a face which omits intermediate rows. The functors
for iterated faces agree strictly. Evaluation remains relative.

Projection to the last row gives relative functors
\[
 r:G_\sigma\longrightarrow I_{\tau_r},\qquad
 s:I_{\tau_r}\longrightarrow G_\sigma,
 \qquad rs=1,\qquad 1\Rightarrow sr.
\]
On $(i,k)$, $r$ has value $u_{ir}k$. Cocartesian arrows give the last
transformation. It descends through the endpoint quotient because
its components there are identities. Every additionally marked
cocartesian arrow maps to an identity under $r$, so both functors
are indeed relative. If the last interval is empty, its two endpoints
coincide and the quotient also identifies $A=B$; the same formulas
apply without choosing two different objects.

\textbf{Contractibility of the three carriers.}
Postcomposition with $r,s$ and their transformation gives categorical
homotopy equivalences on $\mathcal A$ and on every $\mathcal E_t$.
The previous lemma therefore makes
$B\mathcal A(G_\sigma;A,B)$ and
$B\mathcal E_t(G_\sigma;A,B)$ contractible.

For $\mathcal W$ use Lemma~\ref{F:marked-hammocks}, not an
unmarked-vertical assertion. An arbitrary natural transformation
between relative functors induces a natural transformation of their
infinity localizations. To check this directly, functors
$C\times[1]\to D$ are functors $C\to\operatorname{Fun}([1],D)$;
an arrow in the latter functor category is an equivalence exactly
when both endpoint arrows are. The localization universal property
therefore identifies localization of $C\times[1]$ at the copies of
the marking with $LC\times[1]$. If the original transformation is
the identity at $A,B$, its naturality gives a homotopy between the
two induced maps of $\operatorname{Map}(A,B)$. Applying this to
$1\Rightarrow sr$ proves that $r,s$ are inverse on these mapping
spaces; no interior component is being inverted by this argument.

The localization of $I_\tau$ at its negative edges is equivalent to
the ordinal with one edge for each positive edge of $\tau$. Indeed,
functors from the free oriented line are choices of arrows on its
edges; requiring the negative ones to be equivalences contracts each
negative block, with contractible choices of inverse and composite.
This verifies the universal property, including higher choices.
The mapping space between its two endpoints is contractible.
Consequently $B\mathcal W(G_\sigma;A,B)$ is contractible too.

\textbf{The relative carrier argument in every degree.}
Here are the details that a one-square argument does not supply.
Represent a finite sphere, disk, or homotopy diagram by a finite
simplicial complex $K\to N\mathcal A(M;a,b)$ after subdivision and
relative simplicial approximation. Let $P(K)$ be its face poset.
The construction above gives strict covariant diagrams over $P(K)$
of the categories $\mathcal W(G_\sigma)$ and
$\mathcal A(G_\sigma)$. Their classifying spaces are objectwise
contractible. Hence the induced map of their covariant Grothendieck
constructions is a weak equivalence, and each projection to $P(K)$
is a weak equivalence.

The latter construction for $\mathcal A$ has a strict section: at
$\sigma$ choose its last-row path. For a face inclusion
$\theta\subseteq\sigma$, use the universal arrow from the last
row of $\theta$ to the last row of $\sigma$. Composition is the
defining composition of the universal diagram. After evaluation in
$M$ this section is exactly the last-vertex map of the barycentric
subdivision of the original simplex diagram.

If a subcomplex $L\subseteq K$ already maps to
$N\mathcal W(M;a,b)$, every corresponding cocartesian component is
marked by construction. The very same section on $P(L)$ therefore
lifts to the marked Grothendieck construction. Replace its weak
equivalence to the arbitrary-vertical construction by the mapping-path
fibration. It is a trivial fibration. The inclusion
$|P(L)|\subseteq|P(K)|$ is a CW cofibration, so the section lifts
relative to this prescribed boundary. Retracting the mapping-path
replacement yields a lift up to a homotopy fixed on the boundary.
Evaluation in $M$ proves vanishing of every relative homotopy group,
including the path and component assertions. This proves the first
displayed weak equivalence.

Use exactly the same argument with $\mathcal E_t(G_\sigma)$ in
place of $\mathcal W(G_\sigma)$. On a boundary simplex coming from
$\mathcal E_t(M)$ retain its given stage at each row. These stages
decrease along its arrows. For a face inclusion the last-row stages
and the given sequence transition consequently give the required
arrow in $\mathcal E_t(G_\sigma)$; all face compositions still agree.
Thus the boundary section again lifts strictly. Its objectwise
contractible carriers prove that the specified
$B\mathcal E_t(M)\to B\mathcal A(M)$ is a weak equivalence.
No assertion that $t$ is cofinal for arbitrary presheaves is used;
Proposition~\ref{F:C23-uniform} explains why that assertion would
be false.

\textbf{The actual colimit and composition.}
Write $K_n=F_{t_n}(a,b)$, with its actual pullback transitions
$K_m\to K_n$ for $m\leq n$. Sending a stage object to its colimit
class defines the left-hand functor in the stated span: an arrow
$p\to t(n\to m)^*q$ becomes an arrow from $[p]$ to $[q]$.
Any finite diagram in $\mathcal E_t$ pulls back to a common later
stage $N$. Its pullback has a canonical transformation to the original
diagram in the fibred total category, given by identity fibre arrows.
Conversely every finite list of colimit objects, arrows and equations
occurs at a common stage. Applying these two facts to finite sphere
and disk diagrams proves that this specified functor induces all
homotopy-group isomorphisms. It is the fibred version of the filtered
comparison in Proposition~\ref{F:C22-types}.

Finally concatenation is a strict functor on the all-type categories
and commutes with their inclusion $\mathcal W\to\mathcal A$.
The marked comparison identifies it with localized composition by
Lemma~\ref{F:marked-hammocks}. The sequence span therefore transports
the same composition. Where block concatenation gives an actual
stage operation, as for powers of a fixed mixed type, its square with
the all-type functor commutes strictly. Otherwise a chosen later
type receives the concatenated finite data, and the type-map
contraction makes the induced operation agree with this transported
one. Thus the compatibility is with the specified comparisons, not
only an abstract equivalence of individual mapping spaces.
\end{proof}
